\documentclass{article}

\usepackage{graphicx}
\graphicspath{ {figures/} }
\usepackage[hyphens]{url}

% Load last
\usepackage{hyperref}


\bibliographystyle{abbrv}

\begin{document}

%%%%%%%%%%%%%%%%%%%%%%%%%%%%%%%%%%%%%%%%%%%%%%%%%%%%%%%%%%%%%%
\title{Git Branching Strategy}
\author{Matthew Kuperus Heun}
%%%%%%%%%%%%%%%%%%%%%%%%%%%%%%%%%%%%%%%%%%%%%%%%%%%%%%%%%%%%%%

\maketitle


%%%%%%%%%%%%%%%%%%%%%%%%%%%%%%%%%%%%%%%%%%%%%%%%%%%%%%%%%%%%%%
\begin{abstract}
%%%%%%%%%%%%%%%%%%%%%%%%%%%%%%%%%%%%%%%%%%%%%%%%%%%%%%%%%%%%%%

This document gives my strategy for branching in a git repository. 
The strategy is based on the fine blog post by Vincent Driessen
and other resources.
Adherence to this strategy will enable better development processes
\end{abstract}


%%%%%%%%%%%%%%%%%%%%%%%%%%%%%%%%%%%%%%%%%%%%%%%%%%%%%%%%%%%%%%
\section{Introduction} 
\label{sec:introduction}
%%%%%%%%%%%%%%%%%%%%%%%%%%%%%%%%%%%%%%%%%%%%%%%%%%%%%%%%%%%%%%

The use of branches can greatly assist software development projects, 
particularly when developers need to both 
work on new features and maintain released code.

This document gives my model for branching in a git repository. 
The strategy is based on the fine blog post by Vincent Driessen~\cite{Driessen:2010}
and other resources~\cite{Onkelinx:2017, Rankin:2010}.
The model assumes the existence of \texttt{github} or a similar remote repository
and use of the RStudio IDE, 
although the model (if not the details) 
will apply to other IDEs that support branching. 

In addition, it can be difficult to remember 
all of the \texttt{git} commands for branching and merging.
This document gathers all commands in a single place
and discusses the commands in the context of the overall workflow.
 
Adherence to this strategy and use of these commands 
will improve my software development processes.


%%%%%%%%%%%%%%%%%%%%%%%%%%%%%%%%%%%%%%%%%%%%%%%%%%%%%%%%%%%%%%
\section{Workflow} 
\label{sec:workflow}
%%%%%%%%%%%%%%%%%%%%%%%%%%%%%%%%%%%%%%%%%%%%%%%%%%%%%%%%%%%%%%

The workflow presented here assumes familiarity 
with the Driessen branching model. 
All \texttt{git} commands can be entered 
at the command line of the developer's local computer
in the root directory of the repository.
Both Terminal.app (MacOS) and the Terminal window of RStudio
can be used for this purpose.


%++++++++++++++++++++++++++++++
\subsection{Setup} 
\label{sec:setup}
%++++++++++++++++++++++++++++++

Use the following workflow to establish a repository and 
create the master (release) and develop branches.
%
\begin{enumerate}

  \item Create a git repository

  \item Clone the repository, duplicating the \texttt{master} branch on the local computer:
  		\texttt{git clone git@github.com:MatthewHeun/<<RepositoryName>>.git}
  
  \item Create the \texttt{develop} branch: 
  		\texttt{git checkout -b develop master}
		
  \item Make an initial push of the \texttt{develop} branch to origin (github):
        \texttt{git push -u origin develop}
		
  \item Future pushes to github from the \\texttt{develop} branch 
  		can be made from RStudio's \texttt{git interface}.

  \item If needed, click the refresh icon 
  		(\raisebox{-0.3\height}{\includegraphics[scale=0.5]{refresh}}) in the Git pane of RStudio to 
  		show that the project is now on the \texttt{develop} branch.
  
\end{enumerate}


%++++++++++++++++++++++++++++++
\subsection{Develop a feature (\texttt{newfeature})} 
\label{sec:feature}
%++++++++++++++++++++++++++++++

\begin{enumerate}

  \item Create a feature branch: \texttt{git checkout -b newfeature develop} 

  \item If this feature is to be pushed to github, RStudio won't be able to do so until you
  %
  \begin{enumerate}

    \item Commit changes

    \item Make an initial push of the \texttt{newfeature} branch
	      to the origin (github):\\
		  \texttt{git push -u origin newfeature} 
	
	\item Future pushes to github can be made using the RStudio interface.

  \end{enumerate}
  %
  \item When complete, merge \texttt{newfeature} with the \texttt{develop} branch:
  %
  \begin{enumerate}

    \item Switch to develop branch: \texttt{git checkout develop}

    \item Merge \texttt{newfeature} into develop: \texttt{git merge --no-ff newfeature}

  \end{enumerate}
  %
\end{enumerate}


%++++++++++++++++++++++++++++++
\subsection{Prepare for release} 
\label{sec:prep-for-release}
%++++++++++++++++++++++++++++++

\begin{enumerate}

  \item Switch to develop branch: \texttt{git checkout develop}

  \item Create a release branch: \texttt{git checkout -b release-x.x.x develop} 

  \item Update documentation to reflect new release number (\texttt{x.x.x}) and date.
  
  \item Prepare release notes.
  
  \item Commit locally.
  
  \item Push to origin: 

\end{enumerate}
%






\section{Conclusion}

Write your conclusion here.




\bibliography{references}



\end{document}
